% Slide 4 – map of the live demo (brief, section 6: define the group's needs, check a place, show concrete
% barriers and facilities, the source and date of each fact, how incomplete and sample data are marked).
% The slide stays on screen just before switching to the app; details are in the notes.
\begin{frame}{Demo: Karolina plans a trip to a café}
  \begin{columns}[T,onlytextwidth]
    \begin{column}{0.63\textwidth}
      \small
      Karolina is planning to go to a café with someone close to her who uses a wheelchair.

      \vspace{0.5mm}
      \begin{enumerate}
        \setlength{\itemsep}{0.4mm}
        \item \textbf{She sets the needs}\\
              Wheelchair + step-free entrance + accessible toilet
        \item \textbf{Accessly suggests places}\\
              A ranking of cafés matched to the profile, showing:
              \begin{itemize}
                \setlength{\itemsep}{0pt}\setlength{\topsep}{0pt}
                \item what is accessible,
                \item what is missing,
                \item how up to date the data is.
              \end{itemize}
        \item \textbf{She checks the details}\\
              Next to every fact she sees: source \textperiodcentered\ date \textperiodcentered\ verification status
        \item \textbf{She decides}\\
              She picks a place she can trust -- before they leave home.
      \end{enumerate}
    \end{column}
    \begin{column}{0.34\textwidth}
      \screenshot{karta-dopasowanie}\par
      \vspace{0.3mm}
      \muted{Step 2: what the place meets, what is missing, what is unknown.}
      \vspace{1.5mm}
      \screenshot{karta-atrybut}\par
      \vspace{0.3mm}
      \muted{Step 3: value, source (OpenStreetMap), date (2~Oct 2026) and verification status.}
    \end{column}
  \end{columns}

  \vfill
  {\centering\color{accNavy}\bfseries
    Accessly helps carers plan an outing without guesswork or stress on arrival.\par}
  \vspace{1mm}

  \note[item]{Time: about 3 minutes. Before going on stage: the app open on a phone and a laptop, the “Wheelchair” profile selected, the map on the Main Square (zoom 16), the owner panel open in a second browser tab.}
  \note[item]{Step 1 (0:00–0:30): tap the “Wheelchair” chip; show the toast with the report count and the “Nearest to the map centre” list.}
  \note[item]{Step 2 (0:30–1:00): Places tab, type “A café in the centre I can enter in a wheelchair, ideally with an accessible toilet”; in keyword mode the result is deterministic. Read out what it meets and what is “No data”.}
  \note[item]{Step 3 (1:00–1:45): open the card; point at three things next to the attribute: source, date, status. Then “No data”: our most important message, we do not guess. Buttons Up to date / Outdated / Correct, “Add” and a question to the owner.}
  \note[item]{Step 4 (1:45–2:20): “Report a problem”, pin at the stop, type “No lift”, severity “Blocks”; show the marker shape (square = blocks) and the votes. A report needs no account.}
  \note[item]{Step 5 (2:20–3:00): in the second tab, \texttt{/owner.html} as \texttt{kasia@accessly.test} (Camelot Cafe), click “Potwierdź informacje” (the owner panel is in Polish, for venue staff; say what the button does: confirm the information); back on the card the status has changed.}
  \note[item]{Edge cases (if the jury asks or time is left): attribute history with conflicting votes; the “Data sources” page with snapshot dates; a place with a missing attribute does not get the “Matches your needs” badge.}
  \note[item]{If the network fails: everything except the map tiles works from local data. Sample data: three demo accounts and test reports; the admin view without sign-in carries a “Demo mode: sample data” label.}
\end{frame}
